% !TEX root = ../main.tex
\section{Marco teórico y ejercicios resueltos}

\subsection{Comparativa formal de complejidades teóricas}
La Tabla \ref{tab:complejidad_teorica} resume las cotas asintóticas para las operaciones fundamentales entre un montículo binario estándar y un montículo de Fibonacci.

\begin{table}[H]
    \caption{Comparativa de complejidades temporales entre estructuras de prioridad}
    \label{tab:complejidad_teorica}
    \centering
    \small
    \begin{tabular*}{\textwidth}{@{\extracolsep{\fill}}lcc}
        \toprule
        \textbf{Operación} & \textbf{Montículo Binario (\texttt{heapq})} & \textbf{Montículo de Fibonacci} \\
        \midrule
        \texttt{insert} & $O(\log n)$ (peor caso) & $O(1)$ (amortizado) \\
        \texttt{minimum} & $O(1)$ (peor caso) & $O(1)$ (peor caso) \\
        \texttt{extract-min} & $O(\log n)$ (peor caso) & $O(\log n)$ (amortizado) \\
        \texttt{decrease-key} & $O(\log n)$ (con búsqueda) & $O(1)$ (amortizado) \\
        \texttt{union} & $O(n)$ o no constante & $O(1)$ (peor caso) \\
        \texttt{delete} & $O(\log n)$ & $O(\log n)$ (amortizado) \\
        \bottomrule
    \end{tabular*}
\end{table}

% -------------------------------------------------------------------------
\subsection{Ejercicio resuelto 1: instrumentación de la función potencial}
El primer ejercicio resuelto de la guía implementa la contabilidad del potencial y el cálculo del costo amortizado sobre una operación abstracta.

\begin{lstlisting}[language=Python, caption={Código del Ejercicio Resuelto 1 (función potencial y costo amortizado)}, label={lst:resuelto_1}]
from dataclasses import dataclass

@dataclass(frozen=True)
class EstadoPotencial:
    arboles: int
    marcados: int

    def potencial(self) -> int:
        return self.arboles + 2 * self.marcados

def costo_amortizado(costo_real, antes: EstadoPotencial, despues: EstadoPotencial):
    delta = despues.potencial() - antes.potencial()
    return costo_real + delta, delta

# Caso de prueba oficial de la guia 04
antes = EstadoPotencial(arboles=4, marcados=3)
despues = EstadoPotencial(arboles=7, marcados=0)

amortizado, delta = costo_amortizado(4, antes, despues)

print("Phi antes:", antes.potencial())       # 10
print("Phi despues:", despues.potencial())   # 7
print("Delta Phi:", delta)                   # -3
print("Costo amortizado:", amortizado)       # 1

assert antes.potencial() == 10
assert despues.potencial() == 7
assert delta == -3
assert amortizado == 1
\end{lstlisting}

\begin{lstlisting}[style=consolestyle, caption={Salida de consola del Ejercicio Resuelto 1}]
Phi antes: 10
Phi despues: 7
Delta Phi: -3
Costo amortizado: 1
\end{lstlisting}

\subsubsection*{Análisis técnico}
La operación requirió un costo real de $c_i = 4$ unidades. Sin embargo, al desmarcar nodos y liberar árboles, el potencial se redujo en $\Delta \Phi = -3$. El costo amortizado contable resultante es:
\begin{equation}
    \hat{c}_i = c_i + \Delta \Phi = 4 + (7 - 10) = 1
\end{equation}
Este resultado demuestra que operaciones ocasionalmente costosas pueden tener un costo amortizado de tan solo 1 unidad cuando son sufragadas por el potencial previamente almacenado.

% -------------------------------------------------------------------------
\subsection{Ejercicio resuelto 2: consolidación programada por grados}
El segundo programa modela el comportamiento de las raíces durante la consolidación, enlazando árboles de igual grado mediante un diccionario de grados.

\begin{lstlisting}[language=Python, caption={Código del Ejercicio Resuelto 2 (consolidación simplificada)}, label={lst:resuelto_2}]
from dataclasses import dataclass, field

@dataclass
class Raiz:
    clave: int
    grado: int = 0
    hijos: list = field(default_factory=list)

def enlazar(a, b):
    padre, hijo = (a, b) if a.clave <= b.clave else (b, a)
    padre.hijos.append(hijo)
    padre.grado += 1
    return padre

def consolidar(raices):
    por_grado = {}
    for raiz in raices:
        actual = raiz
        while actual.grado in por_grado:
            otra = por_grado.pop(actual.grado)
            actual = enlazar(actual, otra)
        por_grado[actual.grado] = actual
    return list(por_grado.values())

# Banco de prueba oficial
raices = [
    Raiz(23), Raiz(7), Raiz(21), Raiz(18, 1),
    Raiz(52), Raiz(38, 1), Raiz(17, 1), Raiz(24, 2)
]
resultado = consolidar(raices)
grados = sorted(r.grado for r in resultado)

print([(r.clave, r.grado) for r in resultado])
assert len(grados) == len(set(grados))
assert min(r.clave for r in resultado) == 7
\end{lstlisting}

\begin{lstlisting}[style=consolestyle, caption={Salida de consola del Ejercicio Resuelto 2}]
[(52, 0), (23, 1), (17, 2), (7, 3)]
\end{lstlisting}

\subsubsection*{Observaciones de diseño}
\begin{itemize}
    \item La tabla \texttt{por\_grado} resuelve colisiones de grado en tiempo constante promedio.
    \item Cuando dos raíces colisionan (por ejemplo, grado 1), se enlazan y la raíz resultante adquiere grado 2, lo que puede desencadenar una nueva colisión en cascada.
    \item Al finalizar, todos los árboles resultantes poseen grados estrictamente únicos ($0, 1, 2, 3$), y la clave mínima global ($7$) se conserva en una de las raíces.
\end{itemize}

% -------------------------------------------------------------------------
\subsection{Ejercicio resuelto 3: simulación de decrease-key con heapq}
Para contrastar el montículo de Fibonacci con una biblioteca compilada de alto rendimiento, se modeló una cola de prioridad basada en \texttt{heapq} que soporta actualizaciones perezosas.

\begin{lstlisting}[language=Python, caption={Código del Ejercicio Resuelto 3 (cola de prioridad con heapq perezoso)}, label={lst:resuelto_3}]
import heapq
import itertools

REMOVED = object()

class HeapqPriorityQueue:
    def __init__(self):
        self.heap = []
        self.entries = {}
        self.counter = itertools.count()

    def insert(self, item, priority):
        if item in self.entries:
            self.remove(item)
        entry = [priority, next(self.counter), item]
        self.entries[item] = entry
        heapq.heappush(self.heap, entry)

    def remove(self, item):
        entry = self.entries.pop(item)
        entry[2] = REMOVED

    def decrease_key(self, item, new_priority):
        old_priority = self.entries[item][0]
        if new_priority > old_priority:
            raise ValueError("La nueva clave debe ser menor")
        self.insert(item, new_priority)

    def minimum(self):
        self._discard_removed()
        priority, _, item = self.heap[0]
        return item, priority

    def extract_min(self):
        while self.heap:
            priority, _, item = heapq.heappop(self.heap)
            if item is not REMOVED:
                del self.entries[item]
                return item, priority
        raise IndexError("Cola vacia")

    def _discard_removed(self):
        while self.heap and self.heap[0][2] is REMOVED:
            heapq.heappop(self.heap)

# Prueba oficial de la guia
pq = HeapqPriorityQueue()
pq.insert("a", 20)
pq.insert("b", 12)
pq.insert("c", 30)
pq.decrease_key("c", 5)

assert pq.minimum() == ("c", 5)
assert pq.extract_min() == ("c", 5)
print("Activos:", len(pq.entries))
print("Tamano fisico:", len(pq.heap))
\end{lstlisting}

\begin{lstlisting}[style=consolestyle, caption={Salida de consola del Ejercicio Resuelto 3}]
Activos: 2
Tamano fisico: 3
\end{lstlisting}

\subsubsection*{Lectura e interpretación}
Tras ejecutar la disminución de clave sobre el elemento \texttt{"c"}, la cantidad de elementos lógicamente activos es 2 (\texttt{"a"} y \texttt{"b"}), pero el tamaño físico de la lista de \texttt{heapq} es 3. La entrada obsoleta original \texttt{[30, 2, REMOVED]} permanece almacenada en memoria hasta que alcance la raíz durante una extracción posterior.
